%% 
%% Copyright 2019-2024 Elsevier Ltd
%% 
%% Version 2.4
%% 
%% This file is part of the 'CAS Bundle'.
%% --------------------------------------
%% 
%% It may be distributed under the conditions of the LaTeX Project Public
%% License, either version 1.2 of this license or (at your option) any
%% later version.  The latest version of this license is in
%%    http://www.latex-project.org/lppl.txt
%% and version 1.2 or later is part of all distributions of LaTeX
%% version 1999/12/01 or later.
%% 
%% The list of all files belonging to the 'CAS Bundle' is
%% given in the file `manifest.txt'.
%% 
%% Template article for cas-dc documentclass for 
%% double column output.

%\documentclass[a4paper,fleqn,longmktitle]{cas-dc}
\documentclass[a4paper,fleqn]{cas-dc}

%\usepackage[authoryear,longnamesfirst]{natbib}
%\usepackage[authoryear]{natbib}
\usepackage[numbers]{natbib}

%%%Author definitions
\def\tsc#1{\csdef{#1}{\textsc{\lowercase{#1}}\xspace}}
\tsc{WGM}
\tsc{QE}
\tsc{EP}
\tsc{PMS}
\tsc{BEC}
\tsc{DE}
%%%

\begin{document}
\let\WriteBookmarks\relax
\def\floatpagepagefraction{1}
\def\textpagefraction{.001}
\shorttitle{Leveraging social media news}
\shortauthors{J.K. Krishnan et~al.}

\title [mode = title]{Deep learning–based classification of carbonate microfacies components in thin-section images}                      
\tnotemark[1,2]

\tnotetext[1]{This document is the results of the research
   project funded by the National Science Foundation.}

\tnotetext[2]{The second title footnote which is a longer text matter
   to fill through the whole text width and overflow into
   another line in the footnotes area of the first page.}


\author[1]{Musawer Muradi}[orcid=0000-0001-0000-0000]
%\cormark[1]
%\fnmark[1]
\ead{musavir.muradi@gmail.com}
%\ead[url]{www.jkkrishnan.in}

\credit{Conceptualization of this study, Methodology, Software}

%\address[1]{, Street 129, 1043 NX Amsterdam, The Netherlands}

\author[2]{Arnaud Gallois}
\ead{arnaud.gallois@kocaeli.edu.tr}

\author[1]{Alev Mutlu}%[
   %role=Co-ordinator,
   %suffix=Jr,
   %]
%\fnmark[2]
\ead{alev.mutlu@kocaeli.edu.tr}
%\ead[URL]{https://www.university.org}

\credit{Data curation, Writing - Original draft preparation}

\affiliation[1]{organization={Department of Computer Engineering, Kocaeli University},
                city={Kocaeli},
                country={Turkiye}}

\affiliation[2]{organization={Department of Geological Engineering, Kocaeli University},
                city={Kocaeli},
                country={Turkiye}}

%\author[1,3]{T. Rafeeq}
%\cormark[2]
%\fnmark[1,3]
%\ead{t.rafeeq@example.in}
%\ead[URL]{www.campus.in}

\cortext[cor1]{Corresponding author}
\cortext[cor2]{Principal corresponding author}
\fntext[fn1]{This is the first author footnote, but is common to third
  author as well.}
\fntext[fn2]{Another author footnote, this is a very long footnote and
  it should be a really long footnote. But this footnote is not yet
  sufficiently long enough to make two lines of footnote text.}

\nonumnote{This note has no numbers. In this work we demonstrate $a_b$
  the formation Y\_1 of a new type of polariton on the interface
  between a cuprous oxide slab and a polystyrene micro-sphere placed
  on the slab.
  }

\begin{abstract}
Automated quantitative analysis of carbonate microfacies is severely hindered by extreme class imbalance, where background grains such as peloids often constitute over 85\% of samples while environmentally critical grains represent less than 2\%. Standard deep convolutional neural networks optimized with flat cross-entropy fail categorically under these conditions, acting as majority-class detectors. To overcome this limitation, we introduce a mathematical framework that embeds geological taxonomy directly into neural architecture. We propose a Hierarchical ResNet-18 branching network integrated with Stage-Wise Dynamic Focal Loss and Targeted Staged Sampling. This structure isolates overlapping morphological features and shields localized network heads from majority-class gradient domination. Validated using a rigorous Stratified Group K-Fold strategy to prevent image leakage, our hierarchical methodology achieves a state-of-the-art Balanced Accuracy of 96.2\% on an unseen test set, successfully maintaining near-perfect recall for extreme minority classes and establishing a robust foundation for high-resolution facies analysis.
\end{abstract}

\begin{graphicalabstract}
\includegraphics{figs/cas-grabs.pdf}
\end{graphicalabstract}

\begin{highlights}
\item Hierarchical ResNet-18 encodes carbonate grain taxonomy into stage-wise binary decisions.
\item Dynamic stage-specific focal loss and targeted oversampling mitigate extreme class imbalance.
\item The final model achieves 96.2\% balanced accuracy on a strictly grouped unseen test set.
\end{highlights}

\begin{keywords}
Carbonate Microfacies \sep Deep Learning \sep Class Imbalance \sep Hierarchical Classification \sep Focal Loss
\end{keywords}


\maketitle

\section{Introduction}

Automated petrographic analysis of carbonate rocks represents a critical frontier in computational geosciences. Traditional thin-section analysis is inherently subjective, time-intensive, and requires specialized domain expertise to distinguish between morphologically similar grains such as peloids, ooids, and intraclasts. While the application of deep learning has shown promise in automating this process, existing approaches frequently hit a hard performance ceiling when deployed on real-world geological datasets.

The central computational challenge preventing broader adoption is \textbf{extreme class imbalance}. In typical carbonate samples, structureless background grains (like peloids or matrix) can constitute over 85\% of a dataset, while environmentally significant grains (such as broken ooids indicating high-energy reworking) may represent less than 2\%. When presented with this distribution, standard deep convolutional neural networks (DCNNs) optimized via flat Cross-Entropy (CE) loss fail to learn the minority classes. The gradient updates become overwhelmingly dominated by the majority class, leading to a network that achieves high overall accuracy by effectively acting as a majority-class detector, while rendering the recall of critical rare facies near zero.

To overcome the catastrophic failure of flat architectures on imbalanced sedimentary data, this study proposes a novel synthesis of topological priors and dynamic loss landscapes. Our contribution is a mathematically grounded system that combines a \textbf{Hierarchical Decision Structure} with \textbf{Stage-Wise Dynamic Focal Loss Engineering} and \textbf{Targeted Staged Sampling}. By explicitly mapping the taxonomic branching of carbonate grains into the network architecture, we demonstrate how stripping away dominant noise sequentially allows localized network heads to maintain non-zero gradients for extreme minority classes, resulting in state-of-the-art recall across all facies.

\section{Literature Review \& Domain Challenges}

\subsection{The Bottleneck of Petrographic Data Collection}
A persistent challenge in computational geosciences is the acquisition of high-quality, annotated datasets. Automated thin-section analysis differs fundamentally from general computer vision tasks because ground-truth labeling requires expert sedimentological interpretation. Consequently, available training data is often limited in size. Furthermore, individual physical samples naturally exhibit extreme variations in component abundance. A single thin-section is frequently dominated by a pervasive background matrix or a single grain type, such as unstructured peloids, severely limiting the sampling of critical rare components \citep{koeshidayatullah2020}. This inherent class imbalance has been repeatedly cited as a primary bottleneck for deep learning models deployed in automated core analysis \citep{alsultan2024}.

\subsection{Distinguishing Microfacies for Fluid Flow Prediction}
A substantial portion of existing geological machine learning research focuses on macroscopic rock classification---such as differentiating basic igneous, sedimentary, and metamorphic families---or identifying highly distinct fossil macro-structures. However, high-resolution reservoir characterization, which is critical for predicting porosity, permeability, and tracing fluid flow (e.g., in groundwater aquifer management or hydrocarbon exploration), relies on microfacies classification.

This task is significantly more challenging than macroscopic classification. Carbonate grains frequently exhibit overlapping morphological base structures. For example, distinguishing between a completely micritized ooid and a standard peloid operates on a knife-edge decision boundary, relying on the detection of faint, high-frequency spatial gradients such as residual concentric laminae or jagged breakage planes. Single-stage classifiers often conflate these classes because the overarching geometric features overlap heavily in the shared latent space.

\subsection{Algorithmic Limitations on Imbalanced Data}
Standard deep convolutional neural networks (DCNNs) struggle profoundly when trained on these overlapping, imbalanced geological datasets. When optimized via standard cross-entropy loss, the network overwhelmingly prioritizes the abundant classes. The gradient updates become dominated by the defining features of the majority class, leading to near-zero recall for rare but environmentally significant indicators, such as broken ooids \citep{ferreira2020}. While previous research has attempted to mitigate this via generic data augmentation, balanced undersampling, or synthetic data generation (SMOTE), these methods frequently fail to alter the fundamental gradient descent behavior when physical grain morphologies overlap so heavily.

\subsection{Loss Constriction and Structural Priors}
To combat the gradient dilution caused by majority classes, geological studies have increasingly explored dynamic loss landscapes. For instance, recent architectures (e.g., SENet-ConvNeXt-FL) have integrated Focal Loss to mitigate class imbalance within rock thin-section image classification, successfully down-weighting easily classified background examples \citep{dong}. Concurrently, hierarchical classification schemes have been proposed to organize complex petrographic data into primary and secondary categories. Building upon these foundations, our study extends this logic by applying dynamically scaled Focal Loss explicitly across a multi-stage, branching taxonomic network, isolating overlapping features at each node of the decision tree.

\section{Geological Setting \& Dataset Origins}

\subsection{Origin of Carbonate Samples}
The robust training of a microfacies classification model requires a dataset that genuinely reflects the deep-time heterogeneity of carbonate platforms. The thin sections utilized in this study are derived from marine carbonate successions, which inherently record complex depositional and diagenetic histories. Carbonate systems are predominantly biological in origin, rendering their grains highly susceptible to post-depositional alteration (diagenesis), such as micritization, which frequently obscures primary depositional fabrics. It is this exact biological and chemical degradation that forms the overlapping, ambiguous morphologies (e.g., micritized ooids masquerading as peloids) that our computational framework is engineered to resolve.

\subsection{Sample Preparation and Digitization}
To transition from physical geology to a computable dataset, rock samples were impregnated with blue epoxy resin under vacuum to highlight porosity before being cut and ground to a standard petrographic thickness of 30 $\mu$m. These standard thin sections were subsequently digitized using high-resolution optical microscopy.
For deep learning ingestion, the digitized images were subjected to rigorous sedimentological annotation, resulting in 2,642 distinct grain instances. These isolated grain patches formed the foundational training matrix upon which our model was trained.

\section{Methodology \& Mathematical Framework}

\subsection{Problem Formulation \& Dataset Representation}
Let the carbonate grain dataset be defined mathematically as $\mathcal{D} = \{(x_i, y_i)\}_{i=1}^N$, consisting of $N$ individually annotated grains.
Each input $x_i \in \mathbb{R}^{C \times H \times W}$ represents a masked, grain-centered patch (where $H=W=96$ and $C=3$ for RGB channels). The mask acts as a spatial attention mechanism, enforcing that all background pixels are zeroed, isolating the grain's internal morphology and complex boundaries from the cementing matrix.

The objective is to map each $x_i$ to a categorical label $y_i \in \{Peloid, Ooid, Intraclast, Broken\_Ooid\}$. The fundamental challenge is the severe imbalance in the prior probability distribution $P(Y=y)$, where $P(Peloid) \approx 0.85$ and $P(Broken\_Ooid) \approx 0.01$.

\subsection{Re-framing the Problem: Hierarchical Taxonomic Prior}
Standard deep learning architectures map inputs to a flat softmax output, predicting $P(Y=y | X)$ directly. However, in our heavily imbalanced visual domain with overlapping features (e.g., micritized ooids vs. peloids), this approach collapses.

To overcome this, we decompose the classification into a sequence of binary conditional probabilities that reflect the structural taxonomy of the grains. Instead of a single classification, we predict:
\begin{enumerate}
  \item \textbf{Stage 1:} $P(Peloid | X)$
  \item \textbf{Stage 2:} $P(OoidLike | X, \neg Peloid)$
  \item \textbf{Stage 3:} $P(Broken | X, OoidLike)$
\end{enumerate}

By applying the chain rule of probability, the final class probabilities are constructed logically. For example, the probability of a grain being an \textit{Intraclast} is derived by surviving the first two filters:
\begin{equation}
P(Intraclast) = (1 - P(Peloid)) \times (1 - P(OoidLike))
\end{equation}
This structure prevents the parameters responsible for defining subtle broken edges from being overwhelmed by the gradients generated by thousands of featureless peloids.

\subsection{Feature Extraction Architecture}
We utilize a shared Convolutional Neural Network backbone, specifically ResNet-18 initialized with weights from ImageNet, to extract robust geometric and textural representations.
Let the shared backbone be $f_\theta$, mapping an input $x_i$ to a latent representation $h_i \in \mathbb{R}^d$ (where $d=512$):
\begin{equation}
h_i = f_\theta(x_i)
\end{equation}

This singular latent representation $h_i$ is fed simultaneously into three independent Multi-Layer Perceptrons ($MLP^{(s)}$ for stage $s \in \{1,2,3\}$). Each MLP forms a binary decision head outputting a localized probability:
\begin{equation}
p_i^{(s)} = \sigma(MLP^{(s)}(h_i))
\end{equation}
where $\sigma$ is the sigmoid activation function.

\subsection{Resolving Imbalance: Stage-Wise Dynamic Focal Loss}
Under standard Cross Entropy (CE) loss, the derivative $\frac{\partial \mathcal{L}_{CE}}{\partial \theta}$ is dominated by the majority class. To force the network to learn the minority classes, we implement Focal Loss $(\mathcal{L}_{FL})$:
\begin{equation}
\mathcal{L}_{FL}(p_t) = -\alpha_t (1 - p_t)^\gamma \log(p_t)
\end{equation}

Here, $p_t$ represents the model's assigned probability for the true class.
\begin{itemize}
  \item The \textbf{focusing parameter ($\gamma$)}, set to $2.0$, mathematically reduces the relative loss of well-classified, confident examples (where $p_t$ is high). This forces the network's gradient updates to focus entirely on "hard" examples.
  \item The \textbf{weighting factor ($\alpha_t$)} provides explicit class balancing.
\end{itemize}

A crucial innovation of this methodology is that $\alpha$ is scaled dynamically for each specific hierarchical head to combat local imbalance. Because stage 1 separates the absolute majority (Peloids) from everything else, it requires a different weighting scheme than stage 3, which must force the separation of rare broken ooids.
\begin{itemize}
  \item Stage 1 (Peloid vs Non-peloid): $\alpha=0.25, \gamma=2.0$
  \item Stage 2 (Ooid-like vs Intraclast): $\alpha=0.50, \gamma=2.0$
  \item Stage 3 (Whole vs Broken Ooid): $\alpha=0.75, \gamma=2.0$
\end{itemize}

\subsection{Regularization via Sampling and Inference Augmentation}
To prevent parameter collapse in the deepest, most data-starved head ($MLP^{(3)}$), we employ \textbf{Targeted Staged Sampling} during training. Over-sampling is applied forcefully to the broken ooids, artificially altering the prior probability specifically within the batch calculations for Stage 3 without corrupting the broader feature learning of Stage 1.

During inference, we implement \textbf{Test-Time Augmentation (TTA)} to smooth deterministic decision boundaries. TTA is formulated as returning the mathematical expectation of the model's prediction over a set $\mathcal{T}$ of $K$ valid geometric transformations (rotations and flips):
\begin{equation}
\hat{p_i} = \frac{1}{K} \sum_{k=1}^K p(y | T_k(x_i), \theta)
\end{equation}
This expectation reduces variance for morphologically borderline edges, preventing single-pixel structural anomalies from pivoting deep network activations into the wrong class.

\section{Experimental Setup \& Evaluation Metrics}

\subsection{Cross-Validation Strategy}
A recognized pitfall in applying deep learning to geological data is image-level data leakage, where grains from the same thin-section image appear in both training and test sets. Since grains physically adjacent to one another share indistinguishable cementing matrix properties and lighting conditions, random splitting artificially inflates validation accuracy.

To ensure rigorous validation, the 2,642 annotated grains were separated using a \textbf{Strict Stratified Group K-Fold split} (where the "Group" explicitly refers to the source image). The dataset was partitioned firmly into a $60\% / 20\% / 20\%$ Training/Validation/Test split. The final Test set (529 grains) was completely held out during all phase development and threshold tuning, providing a mathematically robust measure of generalization.

\subsection{Computational Implementation Details}
To ensure strict reproducibility, all hierarchical models and baseline flat architectures were implemented natively in PyTorch. Model training, evaluation, and ablation studies were accelerated using a Google Colab environment equipped with a dedicated NVIDIA GPU, exploiting CUDA for highly parallelized tensor computations.

The networks were optimized using the AdamW (Adaptive Moment Estimation with Weight Decay) algorithm. The base learning rate was controlled at $1 \times 10^{-4}$ to ensure stable gradient updates during the fine-tuning of the pre-trained ImageNet ResNet-18 backbone, with specific localized head warmups briefly reaching $5 \times 10^{-4}$. A weight decay penalty of $1 \times 10^{-4}$ was strictly enforced across all parameter groups to provide rigorous $L_2$ regularization, preventing parameter memorization on the heavily imbalanced dataset. Training was conducted utilizing batch sizes ranging from 16 to 64 depending on the staged sampling phase, balancing VRAM limits and gradient stability. To prevent deep overfitting on the noise of minority classes during Focal Loss optimization, aggressive Early Stopping was employed, terminating the training loop if the validation aggregate loss failed to improve for 5 consecutive epochs.

\subsection{Evaluation Metrics for Imbalanced Domains}
Standard Overall Accuracy is fundamentally misleading for evaluating extremely imbalanced geological datasets. Predicting all samples as "Peloid" yields an $87\%$ accuracy, despite completely failing the scientific objective of identifying critical minority grains. Consequently, we prioritize a rigorous suite of threshold-independent metrics:
\begin{itemize}
  \item \textbf{Precision:} The fraction of relevant instances among the predicted instances, penalizing false positives.
  \item \textbf{Recall (Sensitivity):} The fraction of true instances that were successfully retrieved, vital for rare target detection.
  \item \textbf{F1-Score:} The harmonic mean of precision and recall.
  \item \textbf{Balanced Accuracy:} As the primary comparative metric, Balanced Accuracy is formulated as the arithmetic mean of sensitivity (true positive rate) and specificity (true negative rate) across all classes. This mathematically weighs the correct classification of rare broken ooids (N=6) equally against the ubiquitous peloids (N=460) during macro-averaging:
\end{itemize}

\begin{equation}
Balanced Accuracy = \frac{Sensitivity + Specificity}{2}
\end{equation}

\section{Results \& Ablation Study}

\subsection{The Failure of the Flat Baseline Model}
We established a strict baseline using a flat ResNet-18 architecture trained with standard Cross-Entropy. The baseline immediately demonstrated the vulnerability of modern deep learning to sedimentary class imbalance. While it reached strong overall accuracy metrics, the minority classes (Broken Ooids and Intraclasts) experienced near catastrophic collapse in recall, proving unable to establish any meaningful decision boundary distinct from standard Ooids and Peloids.

Conversely, our proposed model structure immediately solved this gradient dominance.

\subsection{Component Ablation Study}
To mathematically prove the contribution of each engineered component in our pipeline, an ablation study was performed on the strict 529-grain test set. Performance metrics are recorded as we progressively re-introduce components to the baseline architecture.

\begin{table*}[t]
\centering
\caption{Test Set Per-Class Recall Progression}
\resizebox{\textwidth}{!}{%
\begin{tabular}{lccccc}
\toprule
\textbf{Model Iteration} & \textbf{Peloid (N=460)} & \textbf{Ooid (N=42)} & \textbf{Broken (N=6)} & \textbf{Intraclast (N=21)} & \textbf{Balanced Acc.} \\
\midrule
Flat ResNet-18 (Baseline) & 452 (98\%) & 33 (79\%) & 5 (83\%) & 12 (57\%) & 94.9\% \\
Flat EfficientNet-B0 + Attn & 439 (95\%) & 35 (83\%) & 2 (33\%) & 5 (24\%) & 89.4\% \\
XGBoost Ensemble Baseline & 444 (97\%) & 33 (79\%) & 4 (67\%) & 11 (52\%) & 93.0\% \\
\textbf{Hierarchical Focal Loss (Base)} & 447 (97\%) & 35 (83\%) & 5 (83\%) & 10 (48\%) & 94.0\% \\
+ Stage 3 Target Oversample (15x) & 450 (98\%) & 37 (88\%) & 5 (83\%) & 16 (76\%) & 96.0\% \\
+ Heavy Oversample (15x) + TTA & 451 (98\%) & 39 (93\%) & 5 (83\%) & 14 (67\%) & \textbf{96.2\%} \\
\bottomrule
\end{tabular}%
}
\end{table*}

\textit{Note: The baseline model here represents the system before the staged augmentations}.

\subsection{Test-Time Augmentation (TTA) Impact}
Table 1 demonstrates a clear progressive increase in overall classification power. Crucially, the introduction of target oversampling specific to Stage 3 prevents $MLP^{(3)}$ parameter collapse, lifting Intraclast recall from 48\% to 76\%. The application of TTA smoothed deterministic boundaries, lifting Ooid recall to a peak of 93\%.

\subsection{Final Model Confusion Matrix}
The final deployed hierarchical model with full Focal scaling, Targeted Oversampling, and TTA achieved a near-perfect separation matrix, suppressing the noise of the dominant 87\% Peloid class without degrading its own recall.

\begin{table}[h]
\centering
\caption{Confusion Matrix (Hierarchical ResNet-18 + Focal Loss + TTA + Staged Oversample)}
\begin{tabular}{lcccc}
\toprule
\textbf{Actual \textbackslash\ Predicted} & \textbf{Peloid} & \textbf{Ooid} & \textbf{Broken} & \textbf{Intraclast} \\
\midrule
\textbf{Peloid} & 451 & 3 & 0 & 6 \\
\textbf{Ooid} & 2 & 39 & 1 & 0 \\
\textbf{Broken} & 0 & 0 & 5 & 1 \\
\textbf{Intraclast} & 5 & 2 & 0 & 14 \\
\bottomrule
\end{tabular}
\end{table}

This matrix proves the mathematical efficacy of the hierarchical split. By forcing the network through the independent $MLP^{(1)}$ Peloid verification step, the subsequent $MLP^{(2)}$ and $MLP^{(3)}$ heads are shielded from Peloid-dominated gradients, allowing them to perfectly recall 5 out of 6 extreme minority Broken Ooids.

\section{Discussion}

\subsection{Feature Separation in the Latent Space}
A recognized limitation in recent automated petrography is the severe recall degradation caused by the natural imbalance of carbonate facies. The fundamental failure of the baseline flat ResNet-18 is its inability to construct distinct geometric manifolds for minor classes in the latent feature space $h_i$. Because Peloids share macroscopic textural features with heavily micritized Ooids and Intraclasts, the dominant class pulls the minority representations toward its center.

By restructuring the optimization environment into three conditional branches, we mathematically enforce feature separation. $MLP^{(1)}$ becomes a highly specialized texture filter, discarding 87\% of the dataset noise. This taxonomic pump allows $MLP^{(2)}$ and $MLP^{(3)}$ to dedicate their entire parameter space to identifying the fine-grained, high-frequency spatial gradients that define internal layering and fragmented edges, features that were previously washed out in the flat architecture.

\subsection{The Regularizing Power of TTA on Borderline Morphology}
As demonstrated in Table 1, Test-Time Augmentation provided the final mathematical lift required to stabilize Ooid and Intraclast recall. Deep Convolutional layers operate with spatial invariance, but they are not strictly rotation-invariant. In visual setups where grains are rotated arbitrarily during original physical deposition, edge cases frequently fall on the absolute borderline of the network's learned decision boundary. By computing the mathematical expectation $\hat{p_i}$ across a set of diverse geometric states, TTA functionally smooths the hyperspace boundary, drastically reducing false-negative variance for morphologically ambiguous grains.

\section{Conclusion}

This study addresses the critical barrier of extreme class imbalance in the automated quantitative analysis of carbonate microfacies. We demonstrate mathematically and empirically that standard flat deep convolutional neural networks optimized via cross-entropy fail categorically at recalling rare, environmentally critical grains (e.g., Broken Ooids at 1\% distribution).

To overcome this, we formulated a novel classification pipeline that embeds the topological logic of geological taxonomy directly into neural architecture. By constructing a Hierarchical ResNet-18 branching network, applying dynamically scaled Stage-Wise Focal Loss, and implementing explicit data over-sampling at precise nodes of the decision tree, we fundamentally altered the optimization landscape. We successfully shielded minority-class parameters from majority-class gradient domination. Evaluated via a rigorous Stratified Group K-Fold strategy to prevent image leakage, our proposed methodology achieved a Balanced Accuracy of 96.2\% on an unseen test set, establishing a mathematically robust foundation for high-resolution facies analysis in the computational geosciences.

\begin{comment}
The package is available at author resources page at Elsevier
(\url{http://www.elsevier.com/locate/latex}).
The class may be moved or copied to a place, usually,\linebreak
\verb+$TEXMF/tex/latex/elsevier/+, %$%%%%%%%%%%%%%%%%%%%%%%%%%%%%
or a folder which will be read                   
by \LaTeX{} during document compilation.  The \TeX{} file
database needs updation after moving/copying class file.  Usually,
we use commands like \verb+mktexlsr+ or \verb+texhash+ depending
upon the distribution and operating system.

\section{Front matter}

The author names and affiliations could be formatted in two ways:
\begin{enumerate}[(1)]
\item Group the authors per affiliation.
\item Use footnotes to indicate the affiliations.
\end{enumerate}
See the front matter of this document for examples. 
You are recommended to conform your choice to the journal you 
are submitting to.

\section{Bibliography styles}

There are various bibliography styles available. You can select the
style of your choice in the preamble of this document. These styles are
Elsevier styles based on standard styles like Harvard and Vancouver.
Please use Bib\TeX\ to generate your bibliography and include DOIs
whenever available.

Here are two sample references: 
\cite{Fortunato2010}
\cite{Fortunato2010,NewmanGirvan2004}
\cite{Fortunato2010,Vehlowetal2013}

\section{Floats}
{Figures} may be included using the command,\linebreak 
\verb+\includegraphics+ in
combination with or without its several options to further control
graphic. \verb+\includegraphics+ is provided by {graphic[s,x].sty}
which is part of any standard \LaTeX{} distribution.
{graphicx.sty} is loaded by default. \LaTeX{} accepts figures in
the postscript format while pdf\LaTeX{} accepts {*.pdf},
{*.mps} (metapost), {*.jpg} and {*.png} formats. 
pdf\LaTeX{} does not accept graphic files in the postscript format. 

\begin{figure}
	\centering
	\includegraphics[width=.9\columnwidth]{figs/cas-munnar-2024.jpg}
	\caption{The beauty of Munnar, Kerala. (See also Table \protect\ref{tbl1}).}
	\label{FIG:1}
\end{figure}


The \verb+table+ environment is handy for marking up tabular
material. If users want to use {multirow.sty},
{array.sty}, etc., to fine control/enhance the tables, they
are welcome to load any package of their choice and
{cas-dc.cls} will work in combination with all loaded
packages.

\begin{table}[width=.9\linewidth,cols=4,pos=h]
\caption{This is a test caption. This is a test caption. This is a test
caption. This is a test caption. Use \{table*\} instead of \{table\} if you
want a two column spanned table.}\label{tbl1}
\begin{tabular*}{\tblwidth}{@{} LLLL@{} }
\toprule
Col 1 & Col 2 & Col 3 & Col4\\
\midrule
12345 & 12345 & 123 & 12345 \\
12345 & 12345 & 123 & 12345 \\
12345 & 12345 & 123 & 12345 \\
12345 & 12345 & 123 & 12345 \\
12345 & 12345 & 123 & 12345 \\
\bottomrule
\end{tabular*}
\end{table}

\section[Theorem and ...]{Theorem and theorem like environments}

{cas-dc.cls} provides a few shortcuts to format theorems and
theorem-like environments with ease. In all commands the options that
are used with the \verb+\newtheorem+ command will work exactly in the same
manner. {cas-dc.cls} provides three commands to format theorem or
theorem-like environments: 

\begin{verbatim}
 \newtheorem{theorem}{Theorem}
 \newtheorem{lemma}[theorem]{Lemma}
 \newdefinition{rmk}{Remark}
 \newproof{pf}{Proof}
 \newproof{pot}{Proof of Theorem \ref{thm2}}
\end{verbatim}


The \verb+\newtheorem+ command formats a
theorem in \LaTeX's default style with italicized font, bold font
for theorem heading and theorem number at the right hand side of the
theorem heading.  It also optionally accepts an argument which
will be printed as an extra heading in parentheses. 

\begin{verbatim}
  \begin{theorem} 
   For system (8), consensus can be achieved with 
   $\|T_{\omega z}$ ...
     \begin{eqnarray}\label{10}
     ....
     \end{eqnarray}
  \end{theorem}
\end{verbatim}  


\newtheorem{theorem}{Theorem}

\begin{theorem}
For system (8), consensus can be achieved with 
$\|T_{\omega z}$ ...
\begin{eqnarray}\label{10}
....
\end{eqnarray}
\end{theorem}

The \verb+\newdefinition+ command is the same in
all respects as its \verb+\newtheorem+ counterpart except that
the font shape is roman instead of italic.  Both
\verb+\newdefinition+ and \verb+\newtheorem+ commands
automatically define counters for the environments defined.

The \verb+\newproof+ command defines proof environments with
upright font shape.  No counters are defined. 

\begin{figure*}
	\centering
	\includegraphics[width=.9\textwidth]{figs/cas-munnar-2024.jpg}
	\caption{The beauty of Munnar, Kerala. (See also Table \protect\ref{tbl1}).}
	\label{FIG:2}
\end{figure*}


\section[Enumerated ...]{Enumerated and Itemized Lists}
{cas-dc.cls} provides an extended list processing macros
which makes the usage a bit more user friendly than the default
\LaTeX{} list macros.   With an optional argument to the
\verb+\begin{enumerate}+ command, you can change the list counter
type and its attributes.

\begin{verbatim}
 \begin{enumerate}[1.]
 \item The enumerate environment starts with an optional
   argument `1.', so that the item counter will be suffixed
   by a period.
 \item You can use `a)' for alphabetical counter and '(i)' 
  for roman counter.
  \begin{enumerate}[a)]
    \item Another level of list with alphabetical counter.
    \item One more item before we start another.
    \item One more item before we start another.
    \item One more item before we start another.
    \item One more item before we start another.
\end{verbatim}

Further, the enhanced list environment allows one to prefix a
string like `step' to all the item numbers.  

\begin{verbatim}
 \begin{enumerate}[Step 1.]
  \item This is the first step of the example list.
  \item Obviously this is the second step.
  \item The final step to wind up this example.
 \end{enumerate}
\end{verbatim}

\section{Cross-references}
In electronic publications, articles may be internally
hyperlinked. Hyperlinks are generated from proper
cross-references in the article.  For example, the words
\textcolor{black!80}{Fig.~1} will never be more than simple text,
whereas the proper cross-reference \verb+\ref{tiger}+ may be
turned into a hyperlink to the figure itself:
\textcolor{blue}{Fig.~1}.  In the same way,
the words \textcolor{blue}{Ref.~[1]} will fail to turn into a
hyperlink; the proper cross-reference is \verb+\cite{Knuth96}+.
Cross-referencing is possible in \LaTeX{} for sections,
subsections, formulae, figures, tables, and literature
references.

\section{Bibliography}

Two bibliographic style files (\verb+*.bst+) are provided ---
{model1-num-names.bst} and {model2-names.bst} --- the first one can be
used for the numbered scheme. This can also be used for the numbered
with new options of {natbib.sty}. The second one is for the author year
scheme. When  you use model2-names.bst, the citation commands will be
like \verb+\citep+,  \verb+\citet+, \verb+\citealt+ etc. However when
you use model1-num-names.bst, you may use only \verb+\cite+ command.

\verb+thebibliography+ environment.  Each reference is a\linebreak
\verb+\bibitem+ and each \verb+\bibitem+ is identified by a label,
by which it can be cited in the text:

\noindent In connection with cross-referencing and
possible future hyperlinking it is not a good idea to collect
more that one literature item in one \verb+\bibitem+.  The
so-called Harvard or author-year style of referencing is enabled
by the \LaTeX{} package {natbib}. With this package the
literature can be cited as follows:

\begin{enumerate}[\textbullet]
\item Parenthetical: \verb+\citep{WB96}+ produces (Wettig \& Brown, 1996).
\item Textual: \verb+\citet{ESG96}+ produces Elson et al. (1996).
\item An affix and part of a reference:\break
\verb+\citep[e.g.][Ch. 2]{Gea97}+ produces (e.g. Governato et
al., 1997, Ch. 2).
\end{enumerate}

In the numbered scheme of citation, \verb+\cite{<label>}+ is used,
since \verb+\citep+ or \verb+\citet+ has no relevance in the numbered
scheme.  {natbib} package is loaded by {cas-dc} with
\verb+numbers+ as default option.  You can change this to author-year
or harvard scheme by adding option \verb+authoryear+ in the class
loading command.  If you want to use more options of the {natbib}
package, you can do so with the \verb+\biboptions+ command.  For
details of various options of the {natbib} package, please take a
look at the {natbib} documentation, which is part of any standard
\LaTeX{} installation.

\appendix
\section{My Appendix}
Appendix sections are coded under \verb+\appendix+.

\verb+\printcredits+ command is used after appendix sections to list 
author credit taxonomy contribution roles tagged using \verb+\credit+ 
in frontmatter.

\printcredits

%% Loading bibliography style file
%\bibliographystyle{model1-num-names}
\bibliographystyle{cas-model2-names}
\end{comment}
% Loading bibliography database
\bibliography{references}


%\vskip3pt

\bio{}
Author biography without author photo.
Author biography. Author biography. Author biography.
Author biography. Author biography. Author biography.
Author biography. Author biography. Author biography.
Author biography. Author biography. Author biography.
Author biography. Author biography. Author biography.
Author biography. Author biography. Author biography.
Author biography. Author biography. Author biography.
Author biography. Author biography. Author biography.
Author biography. Author biography. Author biography.
\endbio

\bio{figs/cas-pic1}
Author biography with author photo.
Author biography. Author biography. Author biography.
Author biography. Author biography. Author biography.
Author biography. Author biography. Author biography.
Author biography. Author biography. Author biography.
Author biography. Author biography. Author biography.
Author biography. Author biography. Author biography.
Author biography. Author biography. Author biography.
Author biography. Author biography. Author biography.
Author biography. Author biography. Author biography.
\endbio

\bio{figs/cas-pic1}
Author biography with author photo.
Author biography. Author biography. Author biography.
Author biography. Author biography. Author biography.
Author biography. Author biography. Author biography.
Author biography. Author biography. Author biography.
\endbio

\end{document}

